\scanpage{23}

\deleted{\unsure{2°}}) \uline{Cas contraire}. Je suis pessimiste.

2°) Cas $p=0$. Par contre $k[[t]]$ sur $k$,
($k=\mathbb Q$, ou $k=\mathbb C\ldots$). Je présume que
si $\dim A\unsure{\geq|>}0$, $A$ n'est jamais \unsure{for}$^t$ simple
sur $k$ pour sa topologie discrète.

Dans les cas défavorables, je doute qu'il
y ait \unsure{un} ``\unsure{form}$^t$ simple contenant un corps
premier'' [si $k$ n'est pas alg. sur le sien].

$\Omega_k\otimes_k A\longrightarrow\Omega_A$
admet un inverse à gauche ?

\deleted{Supposons $k$ \unsure{le degré de} $\unsure{\Omega_k|\Omega^-_k}$}
\[
k=\mathbb Q\bigl((x_i)_{i\in I}\bigr)
\qquad \deleted{\unsure{k=\mathbb Q}}
\]
\[
A=k[[\unsure{t|u}]],\qquad
\Omega_k\otimes_k A\simeq A^{(I)}
\unsure{=}\deleted{\prod}\inserted{\unsure{\bigoplus_{i\in I}}}
A\,d_{/\mathbb Q}x_i,
\]
on cherche $A\unsure{\mathrel{\overset{D}{\text{---}}}|\xrightarrow{D}}\Omega_k\otimes_k A$, qui soit
donné sur $k$ par $D(x_i)=d_{/\mathbb Q}x_i$.
Or s'il y a un tel $D$, alors $Dt$ est combinaison lin.
finie de certains $d_{/\mathbb Q}x_i$, $i\in J$, et \unsure{remplaçant}
$\mathbb Q$ par $k_0=\mathbb Q\bigl((x_i)_{\unsure{i\in J}}\bigr)$,
on peut supposer que
$dt=0$.
\deleted{$D(\sum a_i t^i)=$}

\mspage{22}
